\paragraph{\texorpdfstring{$S_5\supset S_4\supset A_4\supset V_4$}{S5>S4>A4>V4} 子群层的分歧剖分与中间曲线的显式模型}
沿用定理 \ref{thm:xi-terminal-zm-delta-map-hurwitz-passport-s5} 的次数 \(5\) 覆盖
\(\pi_\delta:E\to\PP^1_\delta\) 与结论 \ref{con:xi-terminal-zm-delta-s5-regular-closure-genus-109} 的正则闭包
\[
L/\QQ(\delta),\qquad \Gal(L/\QQ(\delta))\cong S_5,\qquad \QQ(X)=L,
\]
并在自然五点作用下取一点稳定子 \(H\simeq S_4\le S_5\)，从而识别
\[
\QQ(E)=L^{H},\qquad E\simeq X/H,
\]
并记相应商映射为 \(\pi:X\to E\)。

\begin{conclusion}[\texorpdfstring{$S_4$}{S4}-层分歧剖分与分歧点数 \(18\)]\label{con:xi-terminal-zm-delta-s4-layer-branch-decomposition-18}
设 \(\infty\in\PP^1_\delta\) 为 \(\delta\) 的无穷远分歧值，有限分歧值集合为
\[
\{\delta=4\}\ \cup\ \{\delta\in\overline{\QQ}:\ B_5(\delta)=0\}
\]
（定理 \ref{thm:xi-terminal-zm-delta-map-hurwitz-passport-s5}）。
对任意有限分歧值 \(t\)，其纤维在 \(E\) 上分解为
\[
\delta^{-1}(t)=2P_t+Q_{t,1}+Q_{t,2}+Q_{t,3}
\]
（按零点重数计），其中 \(P_t\) 为 \(\pi_\delta\) 的简单分歧点（分歧指数 \(2\)），其余三点 \(Q_{t,i}\) 在 \(\pi_\delta\) 下不分歧。
令
\[
B:=\sum_{t}\ \sum_{i=1}^{3} Q_{t,i},\qquad \deg(B)=18.
\]
则：
\begin{enumerate}
  \item \(\pi:X\to E\) 为 \(S_4\)-伽罗瓦覆叠，\(\deg(\pi)=|S_4|=24\)；
  \item \(\pi\) 在无穷远纤维的唯一点 \(O\in E\) 上方不分歧；
  \item \(\pi\) 的分歧点集恰为 \(B\) 的支撑：对每个 \(Q\in \mathrm{Supp}(B)\)，其惯性群为 \(S_4\) 中由换位生成的阶 \(2\) 循环群，且
  \[
  \#\pi^{-1}(Q)=12,\qquad e_P(\pi)=2\quad(\forall P\in\pi^{-1}(Q));
  \]
  \item \(\deg(R_\pi)=216\)，并与 \(g(X)=109\)、\(g(E)=1\) 的 Hurwitz 公式严格相符。
\end{enumerate}
\end{conclusion}
\begin{proof}
设 \(v\) 为 \(\QQ(\delta)\) 上的分歧赋值，对应惯性群阶 \(|I_v|\in\{5,2\}\)。
取 \(w\) 为 \(\QQ(E)\) 上位于 \(v\) 之上的赋值，则塔式分歧指数满足
\[
e(L/\QQ(E);w)=\frac{e(L/\QQ(\delta);v)}{e(\QQ(E)/\QQ(\delta);w)}
=\frac{|I_v|}{e(E/\PP^1_\delta;w)}.
\]
当 \(v\) 对应无穷远分歧点时，\(|I_v|=5\) 且 \(\pi_\delta\) 在 \(O\) 处全分歧 \(e(E/\PP^1_\delta;O)=5\)，故 \(e(L/\QQ(E);O)=1\)，从而 \(\pi\) 在 \(O\) 处不分歧。
\par
当 \(v\) 对应有限分歧值 \(t\) 时，\(|I_v|=2\)。
在纤维 \(\delta^{-1}(t)\) 中仅有 \(P_t\) 满足 \(e(E/\PP^1_\delta;P_t)=2\)，故 \(e(L/\QQ(E);P_t)=1\)；
而对其余三点 \(Q_{t,i}\) 有 \(e(E/\PP^1_\delta;Q_{t,i})=1\)，从而 \(e(L/\QQ(E);Q_{t,i})=2\)。
因此 \(\pi\) 的分歧点恰为全部 \(Q_{t,i}\)，总数为 \(6\cdot 3=18\)。
由于 \(\pi\) 为 \(S_4\)-伽罗瓦覆叠且每个分歧点惯性阶为 \(2\)，故分歧点之上点数为 \(24/2=12\)，且每点贡献 \(e-1=1\)，从而 \(\deg(R_\pi)=18\cdot 12=216\)。
代入 Hurwitz 公式
\(
2g(X)-2=24(2g(E)-2)+\deg(R_\pi)
\)
即得到 \(g(X)=109\)，与结论 \ref{con:xi-terminal-zm-delta-s5-regular-closure-genus-109} 一致。证毕。
\end{proof}

\begin{conclusion}[主除子恒等式与 \texorpdfstring{$B$}{B} 的解析切出]\label{con:xi-terminal-zm-delta-branch-complement-principal-divisor}
令
\[
R_{\delta,\mathrm{fin}}:=\sum_{t} P_t
\]
为六个有限分歧点之和（按几何点计一次），\(O\in E\) 为 \(\delta\) 的极点。
则：
\begin{enumerate}
  \item 分歧除子满足线性等价 \(R_{\delta,\mathrm{fin}}\sim 6O\)，从而在椭圆曲线群律下 \(\sum_t P_t=O\)；
  \item 微分的除子满足严格恒等式
  \[
  \operatorname{div}(\dd\delta)=R_{\delta,\mathrm{fin}}-6O;
  \]
  \item 令 \(G:=(\delta-4)B_5(\delta)\in\QQ(E)\)，并定义
  \[
  \Theta:=\frac{G}{(\dd\delta)^2}\in\QQ(E),
  \]
  则 \(\Theta\) 的除子为
  \[
  \operatorname{div}(\Theta)=B-18O.
  \]
  特别地，\(\Theta\) 在 \(B\) 上均为一次零点，且在 \(O\) 处有唯一极点阶 \(18\)；并且 \(B\sim 18O\)、\(\sum_{Q\in \mathrm{Supp}(B)}Q=O\)。
\end{enumerate}
\end{conclusion}
\begin{proof}
由 \(K_E\sim 0\) 与 \(K_{\PP^1}\sim -2(\infty)\)，对 \(\pi_\delta:E\to\PP^1_\delta\) 的 Hurwitz 公式给出
\[
R_\delta\sim 2\cdot \pi_\delta^*(\infty)=10O.
\]
又 \(O\) 处全分歧指数为 \(5\)，故
\(
R_\delta=4O+R_{\delta,\mathrm{fin}}
\)，从而 \(R_{\delta,\mathrm{fin}}\sim 6O\)，并在 \(\mathrm{Pic}^0(E)\cong E\) 下等价于 \(\sum_t P_t=O\)。
\par
取 \(\PP^1\) 的局部坐标 \(t\)，则 \(\operatorname{div}(\dd t)=-2(\infty)\)。
由微分拉回公式
\[
\operatorname{div}(\dd\delta)=\pi_\delta^*\operatorname{div}(\dd t)+R_\delta=-2\pi_\delta^*(\infty)+R_\delta=-10O+(4O+R_{\delta,\mathrm{fin}})=R_{\delta,\mathrm{fin}}-6O.
\]
\par
对每个有限分歧值 \(t\)，有
\(
\operatorname{div}(\delta-t)=2P_t+Q_{t,1}+Q_{t,2}+Q_{t,3}-5O
\)。
对 \(t=4\) 与 \(B_5\) 的五个根相乘得到
\[
\operatorname{div}(G)=2R_{\delta,\mathrm{fin}}+B-30O.
\]
于是
\[
\operatorname{div}(\Theta)=\operatorname{div}(G)-2\operatorname{div}(\dd\delta)
=(2R_{\delta,\mathrm{fin}}+B-30O)-2(R_{\delta,\mathrm{fin}}-6O)=B-18O.
\]
由 \(\operatorname{div}(\Theta)\) 的度为零得 \(B\sim 18O\)，并在群律下得到 \(\sum_{Q\in \mathrm{Supp}(B)}Q=O\)。证毕。
\end{proof}

\begin{conclusion}[\texorpdfstring{$A_4$}{A4} 交叉子群与规范化纤维积：属 \(10\) 的二次模型]\label{con:xi-terminal-zm-delta-a4-fiberproduct-y10-explicit}
令 \(A_5\trianglelefteq S_5\) 为交错群，并记
\[
C_{A_5}:=X/A_5,\qquad C_{A_5}:\ s^2=(\delta-4)B_5(\delta)
\]
（结论 \ref{con:xi-terminal-zm-delta-a5-fixedfield-hyperelliptic-genus2}）。
取一点稳定子 \(H\simeq S_4\) 并令 \(A_4:=A_5\cap H\)。
置
\[
Y_{10}:=X/A_4.
\]
则
\[
Y_{10}\ \simeq\ \mathrm{Norm}\bigl(E\times_{\PP^1_\delta} C_{A_5}\bigr),
\]
并且在 \(E\) 上具有显式二次模型
\[
Y_{10}:\ z^2=(\delta-4)B_5(\delta).
\]
此外，二次覆叠 \(Y_{10}\to E\) 的分歧点集恰为 \(B\)，从而 \(g(Y_{10})=10\)。
\end{conclusion}
\begin{proof}
在 \(S_5\)-伽罗瓦扩张 \(L/\QQ(\delta)\) 中，对任意子群 \(U,V\le S_5\) 有域论恒等式
\(
L^{U\cap V}=L^{U}\cdot L^{V}
\)（合成取于 \(L\) 内）。
取 \(U=A_5\)、\(V=H\)，并注意 \(A_5\cap H=A_4\)，得
\[
L^{A_4}=L^{A_5}\cdot L^{H}=\QQ(C_{A_5})\cdot \QQ(E),
\]
其对应的曲线模型即为纤维积 \(E\times_{\PP^1_\delta} C_{A_5}\) 的规范化。
将 \(C_{A_5}\) 的方程在纤维积中取 \(\delta\) 的拉回，即得到 \(z^2=(\delta-4)B_5(\delta)\) 的二次模型。
\par
由结论 \ref{con:xi-terminal-zm-delta-branch-complement-principal-divisor}，
\(
\operatorname{div}(G)=2R_{\delta,\mathrm{fin}}+B-30O
\)，故 \(G\) 在 \(B\) 上为奇阶零点而在 \(R_{\delta,\mathrm{fin}}\) 上为偶阶零点；
因此二次扩张 \(\QQ(E)(\sqrt{G})/\QQ(E)\) 的分歧点集恰为 \(B\)。
由 Hurwitz 公式
\[
2g(Y_{10})-2=2(2g(E)-2)+\deg(B)=18
\]
得到 \(g(Y_{10})=10\)。证毕。
\end{proof}

\begin{conclusion}[\texorpdfstring{$Y_{10}=X/A_4$}{Y10} 的 Jacobian 的 \texorpdfstring{$1+2+7$}{1+2+7} 维同源分解]\label{con:xi-terminal-zm-delta-y10-jacobian-qisogeny-1-2-7}
沿用结论 \ref{con:xi-terminal-zm-delta-a4-fiberproduct-y10-explicit} 的记号，令
\[
E:=X/S_4,\qquad
Y_{10}:=X/A_4,\qquad
C_{A_5}:=X/A_5,
\]
其中 \(S_4\le S_5\) 为一点稳定子，\(A_4:=A_5\cap S_4\)。
则存在在 \(\QQ\) 上定义的阿贝尔簇 \(B_7\)（\(\dim B_7=7\)），使得在 \(\QQ\)-同源意义下
\[
\mathrm{Jac}(Y_{10})\ \sim_{\QQ}\ E\ \times\ \mathrm{Jac}(C_{A_5})\ \times\ B_7.
\]
并且 \(B_7\) 的同源类型由群代数投影（等价地，由 \(S_5\) 在 \(H^0(X,\Omega_X^1)\) 上的等型分量与 \(A_4\)-不变子空间维数）唯一确定。
\end{conclusion}
\begin{proof}[证明要点]
记 \(G:=S_5\)、\(H:=A_4\)。
令
\(
p_H:=\frac{1}{|H|}\sum_{h\in H}h\in \QQ[G]
\)
为群代数投影，则 \(p_H\) 在 \(H^0(X,\Omega_X^1)\) 上的像同构于 \(H^0(Y_{10},\Omega_{Y_{10}}^1)\)，其维数为 \(g(Y_{10})=10\)（结论 \ref{con:xi-terminal-zm-delta-a4-fiberproduct-y10-explicit}）。
\par
由结论 \ref{con:xi-terminal-zm-delta-s5-holomorphic-oneforms-chevalley-weil}，
\(
H^0(X,\Omega_X^1)\cong
\varepsilon^{\oplus 2}\oplus V_{(4,1)}^{\oplus 1}\oplus\cdots\oplus V_{(2,1,1,1)}^{\oplus 7}
\)
为 \(S_5\)-表示直和。
对任一不可约表示 \(V\)，其 \(H\)-不变向量维数由特征平均给出
\[
\dim(V^{A_4})=\frac1{12}\Bigl(\chi_V(1)+3\chi_V((2,2))+8\chi_V((3))\Bigr),
\]
其中 \((2,2)\) 与 \((3)\) 为 \(S_5\) 的双换位与三循环共轭类。
代入 \(S_5\) 特征标表可知：仅当 \(V\in\{\varepsilon,\,V_{(4,1)},\,V_{(2,1,1,1)}\}\) 时 \(\dim(V^{A_4})=1\)，其余不可约表示皆满足 \(\dim(V^{A_4})=0\)。
\par
因此 \(H^0(Y_{10},\Omega^1)\) 仅由上述三类等型分量贡献，且贡献维数分别为
\[
\dim\bigl(\varepsilon^{\oplus 2}\bigr)^{A_4}=2,\qquad
\dim\bigl(V_{(4,1)}^{\oplus 1}\bigr)^{A_4}=1,\qquad
\dim\bigl(V_{(2,1,1,1)}^{\oplus 7}\bigr)^{A_4}=7,
\]
相加得到 \(10\)。
由群代数中心幂等元切出的 \(\QQ\)-同源分解（与结论 \ref{con:xi-terminal-zm-delta-s5-jacobian-qisogeny-decomposition} 的同一口径），
这三块分别对应于：
\begin{itemize}
  \item \(\varepsilon\) 等型分量的 \(A_4\)-不变部分，同源于 \(\mathrm{Jac}(X/A_5)=\mathrm{Jac}(C_{A_5})\)；
  \item \(V_{(4,1)}\) 等型分量的 \(A_4\)-不变部分，为一维椭圆因子，且与 \(E=X/S_4\) 同源；
  \item \(V_{(2,1,1,1)}\) 等型分量的 \(A_4\)-不变部分，定义为 \(B_7\)。
\end{itemize}
故得到所示分解。证毕。
\end{proof}

\begin{conclusion}[二次覆叠 \texorpdfstring{$Y_{10}\to E$}{Y10->E} 的 Prym 分解与极化型]\label{con:xi-terminal-zm-delta-y10-prym-splitting-polarization-type}
令 \(\pi:Y_{10}\to E\) 为自然二次覆叠（由 \(A_4\lhd S_4\) 给出）。
令
\[
P:=\ker\!\bigl(\mathrm{Nm}_\pi:\mathrm{Jac}(Y_{10})\to \mathrm{Jac}(E)\bigr)^0
\]
为 Prym 的连通核分量，\(\Xi\) 为其自然限制极化。
则在 \(\QQ\)-同源意义下
\[
P\ \sim_{\QQ}\ \mathrm{Jac}(C_{A_5})\ \times\ B_7,
\]
并且当 \(\pi\) 的分歧点数为 \(18\)（结论 \ref{con:xi-terminal-zm-delta-s4-layer-branch-decomposition-18}）时，
Prym 极化型满足
\[
\mathrm{type}(\Xi)=(\underbrace{1,\dots,1}_{8},2).
\]
\end{conclusion}
\begin{proof}[证明要点]
由于 \(A_4\lhd S_4\)，有 \(Y_{10}\to E\) 为分歧二次覆叠，且 \(g(E)=1\)。
对任意分歧二次覆叠，标准 Prym 理论给出 \(\QQ\)-同源分裂
\[
\mathrm{Jac}(Y_{10})\ \sim_{\QQ}\ \mathrm{Jac}(E)\times \mathrm{Prym}(Y_{10}/E),
\]
其中 \(\mathrm{Jac}(E)\) 因子为 \(\pi^*\mathrm{Jac}(E)\) 的同源像。
将该分裂与结论 \ref{con:xi-terminal-zm-delta-y10-jacobian-qisogeny-1-2-7} 比较，便得到
\(
P\sim_{\QQ}\mathrm{Jac}(C_{A_5})\times B_7
\)。
\par
极化型由 Mumford 关于分歧二次覆叠的 Prym 极化一般公式给出：若底曲线属为 \(1\)、分歧点数为 \(r\)，则
\(
\mathrm{type}(\Xi)=(1^{r/2-1},2)
\)。
在此 \(r=18\)，故得到 \((1^8,2)\)。证毕。
\end{proof}

\begin{conclusion}[属 \(10\) 曲线 \texorpdfstring{$Y_{10}$}{Y10} 的可割度与低次数态射的分解刚性]\label{con:xi-terminal-zm-delta-y10-gonality-four-factorization}
\[
\operatorname{gon}(Y_{10})=4.
\]
更强地：对任意非恒等态射 \(f:Y_{10}\to\PP^1\)，若 \(\deg f<6\)，则必有 \(\deg f=4\)，并且 \(f\) 必经由二次覆叠 \(\pi:Y_{10}\to E\) 分解，即存在 \(g:E\to\PP^1\) 使得 \(f=g\circ\pi\)。
\end{conclusion}
\begin{proof}[证明要点]
由椭圆曲线到 \(\PP^1\) 的次数 \(2\) 态射 \(E\to\PP^1\) 与 \(\pi\) 复合得到次数 \(4\) 态射 \(Y_{10}\to\PP^1\)，故 \(\operatorname{gon}(Y_{10})\le 4\)。
\par
若存在次数 \(d\le3\) 的态射 \(f:Y_{10}\to\PP^1\)，则 \(f\) 与 \(\pi\) 给出到属分别为 \(0\) 与 \(1\) 的两条曲线的两个独立低次数态射。
由 Castelnuovo--Severi 不等式，
\[
g(Y_{10})\le d\cdot g(E)+2\cdot g(\PP^1)+(d-1)(2-1)=2d-1\le5,
\]
与 \(g(Y_{10})=10\)（结论 \ref{con:xi-terminal-zm-delta-a4-fiberproduct-y10-explicit}）矛盾。
故不存在 \(d\le 3\)，从而 \(\operatorname{gon}(Y_{10})\ge 4\)，结合上界得 \(\operatorname{gon}(Y_{10})=4\)。
\par
再证分解刚性。
若存在次数 \(4\) 态射 \(f\) 不经由 \(\pi\) 分解，则同一不等式给出
\[
g(Y_{10})\le 4\cdot g(E)+(4-1)(2-1)=7,
\]
矛盾。
若存在次数 \(5\) 态射 \(f\)，则 \(f\) 不可能经由次数 \(2\) 的 \(\pi\) 分解；于是
\[
g(Y_{10})\le 5\cdot g(E)+(5-1)(2-1)=9,
\]
亦矛盾。
因此 \(\deg f<6\) 时唯有 \(\deg f=4\) 且必经由 \(\pi\) 分解。证毕。
\end{proof}

\begin{conclusion}[\texorpdfstring{$V_4$}{V4}-resolvent：\texorpdfstring{$S_3$}{S3} 伽罗瓦覆叠与三次无分歧循环上覆]\label{con:xi-terminal-zm-delta-v4-s3-resolvent-z28-etale-c3}
令 \(V_4\trianglelefteq H\simeq S_4\) 为 Klein 四元子群，并置
\[
Z_{28}:=X/V_4.
\]
则 \(Z_{28}\to E\) 为 \(S_3\)-伽罗瓦覆叠（由短正合列 \(1\to V_4\to S_4\to S_3\to 1\) 给出），次数为 \(6\)，且其分歧点集恰为 \(B\)，每个分歧点的惯性群在 \(S_3\) 中为换位生成的阶 \(2\) 子群。
设 \(A_3\trianglelefteq S_3\) 为唯一指数 \(2\) 正规子群，则其逆像为 \(A_4\)，从而 \(Z_{28}\) 的唯一二次中间商为 \(Y_{10}=X/A_4\)，并且自然映射
\[
Z_{28}\longrightarrow Y_{10}
\]
为次数 \(3\) 的无分歧循环覆叠（伽罗瓦群 \(C_3\)）。
\end{conclusion}
\begin{proof}
由 \(V_4\trianglelefteq H\) 且 \(H/V_4\simeq S_3\) 可知 \(Z_{28}\to E\) 为 \(S_3\)-伽罗瓦覆叠且次数为 \(6\)。
在结论 \ref{con:xi-terminal-zm-delta-s4-layer-branch-decomposition-18} 中，\(\pi:X\to E\) 的惯性在 \(H\) 内为换位；
其在商群 \(H/V_4\simeq S_3\) 中仍为换位，故 \(Z_{28}\to E\) 的分歧点集与 \(\pi\) 相同，均为 \(B\)。
\par
在 \(S_3\) 中唯一指数 \(2\) 正规子群为 \(A_3\)；
其在 \(H\) 中的逆像为 \(A_4\)，从而 \(Z_{28}/A_3\simeq X/A_4=Y_{10}\)。
由于 \(A_3\) 为 \(S_3\) 的换位换位积生成的 \(3\)-循环子群，其不含阶 \(2\) 元素，故在 \(B\) 上惯性在该层必为平凡，从而 \(Z_{28}\to Y_{10}\) 无分歧且为循环三次覆叠。证毕。
\end{proof}

\begin{conclusion}[\texorpdfstring{$Z_{28}=X/V_4$}{Z28} 的 Jacobian 的 \texorpdfstring{$1+2+7+6+12$}{1+2+7+6+12} 维同源分解]\label{con:xi-terminal-zm-delta-z28-jacobian-qisogeny-1-2-7-6-12}
沿用结论 \ref{con:xi-terminal-zm-delta-v4-s3-resolvent-z28-etale-c3} 的记号。
则存在在 \(\QQ\) 上定义的阿贝尔簇 \(B_3,B_6\)（分别满足 \(\dim B_3=3\)、\(\dim B_6=6\)），使得在 \(\QQ\)-同源意义下
\[
\mathrm{Jac}(Z_{28})
\ \sim_{\QQ}\
E\ \times\ \mathrm{Jac}(C_{A_5})\ \times\ B_7\ \times\ B_3^{\,2}\ \times\ B_6^{\,2}.
\]
\end{conclusion}
\begin{proof}[证明要点]
令 \(H:=V_4\subset A_4\subset S_5\)，其中 \(H\) 的非平凡元素均为双换位型 \((2,2)\)。
对 \(S_5\) 的不可约表示 \(V\)（特征标 \(\chi_V\)），由特征平均得
\[
\dim(V^{V_4})=\frac14\bigl(\chi_V(1)+3\chi_V((2,2))\bigr).
\]
结合结论 \ref{con:xi-terminal-zm-delta-s5-holomorphic-oneforms-chevalley-weil} 的 Chevalley--Weil 分解，
代入 \(S_5\) 特征标表可得：\(\varepsilon,V_{(4,1)},V_{(2,1,1,1)}\) 的不变维数为 \(1\)，
而 \(V_{(3,2)}\) 与 \(V_{(2,2,1)}\) 的不变维数为 \(2\)；
其余不可约表示的不变维数为 \(0\)。
\par
因此 \(H^0(Z_{28},\Omega^1)=H^0(X,\Omega^1)^{V_4}\) 的贡献维数为
\[
2\cdot 1\ +\ 1\cdot 1\ +\ 7\cdot 1\ +\ 3\cdot 2\ +\ 6\cdot 2
=28,
\]
与 \(g(Z_{28})=28\)（结论 \ref{con:xi-terminal-zm-delta-s5-etale-solvable-tower}）一致。
按群代数中心幂等元的同源分解口径，上述五类贡献在 \(\mathrm{Jac}(Z_{28})\) 上切出五个 \(\QQ\)-同源因子；
其中前三个分别与 \(E\)、\(\mathrm{Jac}(C_{A_5})\)、\(B_7\) 一致，后两类的等型分量分别定义为 \(B_3^{2}\) 与 \(B_6^{2}\)（其中 \(\dim B_3=3,\dim B_6=6\)）。
故得到所示分解。证毕。
\end{proof}

\begin{conclusion}[无分歧三次覆叠 \texorpdfstring{$Z_{28}\to Y_{10}$}{Z28->Y10} 的 Prym：极化型与 \texorpdfstring{$B_3,B_6$}{B3,B6} 的精确识别]\label{con:xi-terminal-zm-delta-z28-y10-prym-polarization-type-b3-b6}
令 \(\pi:Z_{28}\to Y_{10}\) 为结论 \ref{con:xi-terminal-zm-delta-v4-s3-resolvent-z28-etale-c3} 的无分歧循环三次覆叠。
记
\[
P_{28/10}:=\Prym(Z_{28}\to Y_{10}).
\]
则：
\[
\dim P_{28/10}=18,
\qquad
\mathrm{type}(\Xi)=\bigl(1^{9},3^{9}\bigr),
\]
其中 \(\Xi\) 为 Prym 的自然限制极化。
并且在 \(\QQ\)-同源意义下
\[
P_{28/10}\ \sim_{\QQ}\ B_3^{\,2}\times B_6^{\,2}.
\]
\end{conclusion}
\begin{proof}[证明要点]
对无分歧循环素数度覆叠（此处 \(\deg(\pi)=3\)），Prym 的维数与极化型为标准结论：
\[
\dim P_{28/10}=(3-1)\bigl(g(Y_{10})-1\bigr)=18,
\qquad
\mathrm{type}(\Xi)=\bigl(1^{g(Y_{10})-1},3^{g(Y_{10})-1}\bigr)=\bigl(1^{9},3^{9}\bigr),
\]
其中 \(g(Y_{10})=10\) 见结论 \ref{con:xi-terminal-zm-delta-a4-fiberproduct-y10-explicit}。
\par
另一方面，结论 \ref{con:xi-terminal-zm-delta-y10-jacobian-qisogeny-1-2-7} 与 \ref{con:xi-terminal-zm-delta-z28-jacobian-qisogeny-1-2-7-6-12} 给出
\[
\Jac(Y_{10})\sim_{\QQ}E\times \Jac(C_{A_5})\times B_7,
\qquad
\Jac(Z_{28})\sim_{\QQ}E\times \Jac(C_{A_5})\times B_7\times B_3^{\,2}\times B_6^{\,2}.
\]
对无分歧循环覆叠，\(\pi^*\Jac(Y_{10})\) 与 \(\Jac(Z_{28})\) 的不变分量同源，而 Prym 为范数映射的连通核，与该不变分量在 \(\QQ\)-同源意义下互补。
由维数对比（\(28=10+18\)）及上式分解，互补部分恰为 \(B_3^{\,2}\times B_6^{\,2}\)，故得到同源识别。证毕。
\end{proof}

\begin{conclusion}[\texorpdfstring{$A_4$}{A4}-Hecke 代数中的三组正交投影与 \texorpdfstring{$Y_{10}$}{Y10} 的可计算分解]\label{con:xi-terminal-zm-delta-y10-hecke-idempotents-computable-projections}
令 \(G:=S_5\)、\(H:=A_4\)，并令 \(p_H:=\frac1{12}\sum_{h\in H}h\in\QQ[G]\)。
对 \(G\) 的共轭类 \(\mathfrak c\)，令 \(K_{\mathfrak c}:=\sum_{g\in \mathfrak c}g\in Z(\QQ[G])\) 为其类和。
定义三个中心原始幂等元
\[
e_{\varepsilon}=\frac1{120}\Bigl(K_{1}-K_{(2)}+K_{(2,2)}+K_{(3)}-K_{(3,2)}-K_{(4)}+K_{(5)}\Bigr),
\]
\[
e_{\mathrm{std}}=\frac1{30}\Bigl(4K_{1}+2K_{(2)}+K_{(3)}-K_{(3,2)}-K_{(5)}\Bigr),
\]
\[
e_{\varepsilon\otimes\mathrm{std}}=\frac1{30}\Bigl(4K_{1}-2K_{(2)}+K_{(3)}+K_{(3,2)}-K_{(5)}\Bigr).
\]
则 \(e_{\varepsilon},e_{\mathrm{std}},e_{\varepsilon\otimes\mathrm{std}}\) 分别对应于 \(S_5\) 的三种不可约表示
\(\varepsilon\)、\(V_{(4,1)}\) 与 \(V_{(2,1,1,1)}\)。
并且在 Hecke 代数 \(p_H\QQ[G]p_H\) 中，三元
\[
\pi_{\varepsilon}:=p_He_{\varepsilon}p_H,\qquad
\pi_{\mathrm{std}}:=p_He_{\mathrm{std}}p_H,\qquad
\pi_{\varepsilon\otimes\mathrm{std}}:=p_He_{\varepsilon\otimes\mathrm{std}}p_H
\]
两两正交且和等于 \(p_H\)。
因此它们在 \(\mathrm{Jac}(Y_{10})\) 上诱导出三项同源分解
\[
\mathrm{Jac}(Y_{10})
\ \sim_{\QQ}\
\mathrm{Im}(\pi_{\mathrm{std}})\ \times\ \mathrm{Im}(\pi_{\varepsilon})\ \times\ \mathrm{Im}(\pi_{\varepsilon\otimes\mathrm{std}}),
\]
并分别同源于 \(E\)、\(\mathrm{Jac}(C_{A_5})\)、\(B_7\)。
\end{conclusion}
\begin{proof}[证明要点]
中心原始幂等元由一般公式
\(
e_{\lambda}=\frac{\dim V_\lambda}{|G|}\sum_{g\in G}\chi_\lambda(g^{-1})\,g
\)
给出；对对称群特征标取实值，按共轭类整理即得到所示显式表达。
又由结论 \ref{con:xi-terminal-zm-delta-y10-jacobian-qisogeny-1-2-7} 的特征平均计算可知：
除 \(\varepsilon\)、\(V_{(4,1)}\)、\(V_{(2,1,1,1)}\) 外，其余不可约表示在 \(A_4\) 下不含不变向量，
故 \(p_H\) 在 \(\QQ[G]\) 中的中心分解仅保留这三块。
因此 \(p_H e_\lambda p_H\) 在 \(p_H\QQ[G]p_H\) 中互相正交且求和为 \(p_H\)，并在 Jacobian 上诱导同源分裂；同源因子的几何识别由 \(V_{(4,1)}\) 的椭圆商实现（\(E=X/S_4\)）、由 \(\varepsilon\) 的 \(A_5\)-商实现（\(C_{A_5}=X/A_5\)），以及由 \(V_{(2,1,1,1)}\) 的余项定义得到。证毕。
\end{proof}

\begin{conclusion}[三条共轭三次子覆叠与 \texorpdfstring{$S_3$}{S3} 闭包]\label{con:xi-terminal-zm-delta-three-cubic-subcovers-s3-closure}
在 \(S_3\) 中换位子群共有三条共轭类。
取 \(D_8\le H\simeq S_4\) 为 \(V_4\) 的一个上覆，使得 \(D_8/V_4\) 对应于 \(S_3\) 的某一换位子群（从而 \(|D_8|=8\) 且 \([H:D_8]=3\)）。
置
\[
C_{10}:=X/D_8.
\]
则：
\begin{enumerate}
  \item \(C_{10}\to E\) 为次数 \(3\) 的非伽罗瓦覆叠，在 \(B\) 上处处为简单分歧（局部分歧指数 \(2\)），且无其他分歧点；因此 \(g(C_{10})=10\)；
  \item 该覆叠的伽罗瓦闭包为 \(Z_{28}\to E\)；等价地，\(Z_{28}\) 为 \(C_{10}\to E\) 的 \(S_3\)-伽罗瓦闭包；
  \item 在 \(E\) 上存在恰好三条此类次数 \(3\) 覆叠，它们在 \(Z_{28}\) 的 \(S_3\)-作用下相互置换。
\end{enumerate}
\end{conclusion}
\begin{proof}
由 \([H:D_8]=3\) 得 \(C_{10}\to E\) 的次数为 \(3\)。
在 \(B\) 上，\(H\) 的惯性为换位，其在陪集集 \(H/D_8\) 上诱导的置换型为 \(2+1\)，从而 \(C_{10}\to E\) 在 \(B\) 上简单分歧且无额外分歧点。
Hurwitz 公式给出
\[
2g(C_{10})-2=3(2g(E)-2)+\deg(B)=18,
\]
故 \(g(C_{10})=10\)。
\par
由于 \(D_8/V_4\) 为 \(S_3\) 的换位子群，其共轭闭包生成 \(S_3\)；
在 \(H\) 中相应的共轭闭包生成 \(V_4\) 之上全体 \(H\)，故函数域意义下的伽罗瓦闭包为 \(X/V_4=Z_{28}\)。
换位子群在 \(S_3\) 中恰有三条共轭类，故得到三条互为共轭的次数 \(3\) 子覆叠，并被 \(S_3\) 置换。证毕。
\end{proof}

\begin{conclusion}[标准四维等型分量由五个椭圆商生成且与 \texorpdfstring{$E^4$}{E^4} 同源]\label{con:xi-terminal-zm-delta-jacobian-standard-part-e4}
令 \(H_i\le S_5\)（\(i=1,\dots,5\)）为自然五点作用下的五个点稳定子（两两共轭且同构于 \(S_4\)），并记
\[
E_i:=X/H_i,\qquad \pi_i:X\to E_i.
\]
由 \(S_5\) 在 \(X\) 上的甲板变换，曲线 \(\{E_i\}\) 在 \(\QQ\) 上两两同构，且同构于既定椭圆曲线 \(E\)。
令
\[
\iota_i:\ E_i\simeq J(E_i)\xrightarrow{\ \pi_i^*\ }J(X)
\]
为诱导同态，并令
\[
A_{\mathrm{std}}:=\langle\,\iota_1(E_1),\dots,\iota_5(E_5)\,\rangle\subset J(X)
\]
为其生成的最小阿贝尔子簇。
则 \(A_{\mathrm{std}}\) 的维数为 \(4\)，并存在 \(\QQ\)-同源
\[
E^5/\Delta(E)\ \sim_{\QQ}\ A_{\mathrm{std}},
\qquad
 A_{\mathrm{std}}\ \sim_{\QQ}\ E^4,
\]
其中 \(\Delta(E)\subset E^5\) 为对角嵌入。
并且在结论 \ref{con:xi-terminal-zm-delta-s5-jacobian-qisogeny-decomposition} 的记号下，有
\(
A_{\mathrm{std}}\sim_{\QQ}A_{(4,1)}
\)。
\end{conclusion}
\begin{proof}
取每个 \(E_i\) 上非零正则微分 \(\omega_i\)，则 \(\pi_i^*\omega_i\in H^0(X,\Omega_X^1)\) 非零。
由于 \(H_i\) 在 \(S_5\) 中两两共轭，且 \(S_5\) 作用置换这些商映射，故 \(\{\pi_i^*\omega_i\}_{i=1}^{5}\) 的线性张成空间在 \(S_5\) 下同构于五点置换表示。
另一方面，\(H^0(X,\Omega_X^1)^{S_5}\cong H^0(X/S_5,\Omega^1)=H^0(\PP^1_\delta,\Omega^1)=0\)，故该置换表示中平凡不变向量不出现，于是其非平凡部分为标准表示 \(V_{(4,1)}\)，维数为 \(4\)。
因此由 \(\iota_i\) 生成的阿贝尔子簇 \(A_{\mathrm{std}}\) 的正则微分空间对偶于该 \(4\) 维子表示，从而 \(\dim A_{\mathrm{std}}=4\)，并与 \(A_{(4,1)}\) 等型一致。
\par
定义同态
\[
\Phi:\ E_1\times\cdots\times E_5\longrightarrow J(X),\qquad
(P_1,\dots,P_5)\longmapsto \sum_{i=1}^5 \iota_i(P_i).
\]
其像为 \(A_{\mathrm{std}}\)。
对角子簇 \(\Delta(E)\) 的像在微分层面对应置换表示的平凡方向；由于该方向在 \(H^0(X,\Omega^1)\) 中不出现，故 \(\Delta(E)\subset\ker\Phi\)，从而 \(\Phi\) 下降为同态
\[
\overline{\Phi}:\ E^5/\Delta(E)\longrightarrow A_{\mathrm{std}}.
\]
两端维数同为 \(4\)，且 \(\overline{\Phi}\) 在微分上为同构，故其为同源；再由 \(E^5/\Delta(E)\sim_{\QQ}E^4\) 得结论。证毕。
\end{proof}

\endinput
